\section{PERFORMANCE METRICS AND JOINT DESIGN TRADEOFFS}
\label{sec:metrics}
\label{sec:tradeoffs}
After locating a shared choice, the analysis turns to the reported
communication and sensing outcomes. A reported value gains engineering meaning
from its task, its observation point in the signal chain, and its operating and
validation conditions. The comparison profile connects these elements to the
performance evidence.

The primary synthesis contains 4,779 metric records. Each record is an extraction
unit that retains the reported quantity together with the definition and source
conditions available for it. The denominator therefore represents reported
quantities in their original contexts. Data rate, bit error
rate (BER), and range resolution have the widest study coverage, appearing in
85, 60, and 50 studies, respectively. Since one study may report several
families, these totals provide overlapping views of the evidence.

Beyond these familiar measures, the evidence extends from link and waveform
behavior to outcomes observed at the device, estimator, application, and
implementation levels. This breadth makes semantic alignment the organizing
principle for the analysis that follows.

\subsection{Metric Specific Application of the Comparison Profile}

Applying the comparison profile defined in
Section~\ref{sec:foundations_comparison}, each metric row was coded by source
definition, task, measurement plane, operating point, validation method, and
value provenance. Communication and sensing results form a joint operating
point when the source links them to the same condition set.
Table~\ref{tab:metric_admissibility} summarizes domain coverage and the
analytical role assigned to each record.

\begin{table}[!htbp]
\caption{Metric evidence profile for 4,779 primary records}
\label{tab:metric_admissibility}
\centering
\footnotesize
\setlength{\tabcolsep}{3.5pt}
\renewcommand{\arraystretch}{1.04}
\begin{tabularx}{\columnwidth}{@{}>{\raggedright\arraybackslash}X r r@{}}
\toprule
\multicolumn{3}{@{}l}{\emph{Metric domains}} \\
\addlinespace[1pt]
Domain & Records, $n$ (\%) & Studies, $n$ (\%) \\
\midrule
Sensing & 1,816 (38.0\%) & 203 (98.5\%) \\
Communication & 1,328 (27.8\%) & 194 (94.2\%) \\
Joint & 870 (18.2\%) & 158 (76.7\%) \\
Implementation & 476 (10.0\%) & 64 (31.1\%) \\
Smaller normalized domains & 289 (6.0\%) & 33 (16.0\%) \\
\bottomrule
\end{tabularx}

\vspace{3pt}
\begin{tabularx}{\columnwidth}{@{}>{\raggedright\arraybackslash}p{0.275\columnwidth}
>{\raggedleft\arraybackslash}p{0.180\columnwidth}
>{\raggedright\arraybackslash}X@{}}
\toprule
\multicolumn{3}{@{}l}{\emph{Analytical roles}} \\
\addlinespace[1pt]
Role & Records, $n$ (\%) & Use in the survey \\
\midrule
Bounded relation & 118 (2.5\%) & Conditional comparison across studies \\
Native context & 4,661 (97.5\%) & Mechanism, feasibility, and qualitative
synthesis \\
\bottomrule
\end{tabularx}
\vspace{1pt}
\parbox{\columnwidth}{\footnotesize\emph{Note.} Domain rows and analytical
roles each partition the 4,779 primary records, while study counts overlap. The
governed ledger contains 4,861 rows, including 31 for contextual interpretation
and 51 under conflict governance. Comparison follows the reported definition,
task, measurement plane, conditions, and validation. Electronic Supplement
S-Evidence provides the full records and provenance, including quantities
preserved under the residual ``other'' code.}
\end{table}

The domain distribution shows where metric evidence is concentrated, while the
analytical role determines how each record enters the survey. The following
subsections explain how task, measurement plane, and operating conditions shape
communication, sensing, joint, and implementation metrics.

\subsection{Metrics Across Functions}

\subsubsection{Communication Metrics}

A communication rate gains meaning from what the source counts. Symbol rate
tracks signaling pace, while data rate reflects the payload and overhead
included in that count. Throughput describes delivered service, capacity
provides a model bound, and spectral efficiency relates useful data to its
stated denominator. In a photonics assisted terahertz prototype, net rates
follow communication subcarrier allocation, payload occupancy, and coding
overhead
\cite{OISAC_SCR01054}. An embedded pilot orthogonal frequency division
multiplexing (OFDM) simulation defines spectral efficiency as the data
subcarrier fraction and reports cyclic prefix overhead
separately
\cite{OISAC_SCR00007}.

Reliability shifts attention to the event and its location in the signal chain.
BER measured before and after forward error correction describes two stages of
the link, while packet loss and outage describe broader delivery boundaries. A
long span fiber experiment reports receiver OSNR and BER before correction
during concurrent distributed sensing
\cite{OISAC_SCR01180}. A robotics study reports packet loss over distance and
view angle, display delay over queue size, and formation control under dynamic
motion as separate results \cite{OISAC_SCR00910}. The reported reliability value
is therefore tied to its event, observation point, time scale, and system boundary.

The observation point also determines how signal quality and power are
interpreted. OSNR relates signal power to noise in a stated optical reference
bandwidth, while electrical SNR follows photodetection and digital SNR follows
equalization. Launch and received power likewise refer to the transmitter and
receiver ends of the link. A fiber polarization sensing experiment traces
recovered vibration SNR across an OSNR sweep, connecting the
sensing response to a stated optical condition
\cite{OISAC_SCR00220}.

Distance and latency complete the profile by locating the result in space and
time. Fiber span and wireless separation describe link geometry, while target
range and coverage radius describe the observed or served region. Latency may
include propagation and buffering, signal acquisition and processing, or the
full application loop. Stated endpoints, included stages, and operating
conditions define the scope of each value. The complete profile connects each
communication metric to the system behavior it represents.

\subsubsection{Sensing Metrics}

Sensing metrics first take meaning from the task and quantity they represent.
Range and position estimates, motion measurements, images, and class decisions
refer to distinct observations and references. Resolution describes
separability along a stated dimension, while estimation error relates an
estimate to a reference. Repeatability captures variation across trials, and a
model bound supplies a conditional benchmark. These distinctions establish the
comparison basis.

Representative studies show how this contract works. A fiber reflectometry
system reports breakpoint localization error and range resolution separately
for laboratory and simulated conditions \cite{OISAC_SCR00056}. A visible light
localization experiment distinguishes lateral estimation error from the
smallest demonstrated displacement and summarizes signed displacement with
means and standard deviations across repeated trials
\cite{OISAC_SCR00196}. An optical phased array simulation relates imaging root
mean square error to electrical SNR under a defined scene, waveform, and
beamforming constraint \cite{OISAC_SCR00940}. Each result retains its summary,
geometry, and validation setting.

Detection and classification add dataset and decision context. Their metrics
take meaning from class balance, decision thresholds, background conditions,
and the relation between training and test data. CORE-Lens reports recognition
accuracy and optical camera communication BER across stated class, data rate,
and background settings, with runtime evaluated separately
\cite{OISAC_SCR00833}. VehicleTalk evaluates link performance, distance
estimation error, and vehicle status recognition under separate prototype
conditions \cite{OISAC_SCR01034}. Their application meaning follows from the
attached classes, geometries, and test protocols.

Physical quantity sensing follows the same logic by linking each measurand to a
reference and calibration. Spatial support and acquisition time locate the
observation, while the aggregation rule shapes its summary. The resulting
profile connects the task, reference, geometry, processing, and validation
setting to the reported sensing outcome. Together, these elements define the
conditions for comparison.

\subsubsection{Joint and Implementation Metrics}

Joint metrics connect communication and sensing through a shared operating
point or resource rule. A matched tuple uses the same conditions for both
functions. A shared resource can also link values that retain their own
reported conditions. In an embedded pilot OFDM simulation, BER and mean
absolute pilot error are reported at the same SNR with eight pilots, while
spectral efficiency follows the corresponding data subcarrier fraction
\cite{OISAC_SCR00007}. The pilot error measures deviation of known pilot symbols
in the simulated channel. Together, the three values show how pilot allocation
connects communication reliability, resource efficiency, and the reported
sensing measure.

Optimization creates another form of joint evidence by selecting an operating
point through shared design variables. In a simulated optical phased array
system, precoding matrices and photodiode orientations are selected together.
The design maximizes a sensing contrast metric under a communication quality
threshold and emitter power constraints \cite{OISAC_SCR00940}. Geometry,
constraint thresholds, and metric definitions define the scope of the
solution.

Implementation metrics place these operating points within the platform and
application chain. Receiver bandwidth, converter rate, memory, acquisition
time, and latency show where processing burdens arise. An outdoor photonics
assisted terahertz prototype evaluates receiver pipeline latency in hardware
analysis and communication and ranging performance in a field experiment
\cite{OISAC_SCR01054}. A smartphone camera prototype evaluates recognition and
communication in the laboratory and benchmarks module latency, memory, and
complete cycle timing on the device \cite{OISAC_SCR00833}. Together, these
measures connect joint performance to the processing chain and locate the
effects of shared design choices in system behavior.

\subsection{How Communication and Sensing Performance Interact}

Communication and sensing often share bandwidth, power, waveform components,
hardware, and processing. A design change in any of these elements can move
both functions to a new operating point. More pilot symbols may strengthen a
sensing reference while leaving fewer resources for data. Wider bandwidth may
support both a higher data rate and finer range resolution until a receiver or
platform limit becomes active. The evidence is organized by the shared choice
that produces each pattern.

The synthesis covers 402 relationships from 168 studies. Each relationship is
read within its waveform, geometry, target, observation point, and validation
setting. These details matter because 371 relationships depend on the reported
operating conditions. The following discussion explains what is shared, how
both functions respond, and when the pattern can guide another design.

% FUTURE FIGURE 6 --- CONDITION-AWARE TRADEOFF EVIDENCE PROFILE
% Do not generate yet.
% Stable label: fig:tradeoff_profile
% Blueprint ID: FIG-OISAC-05
% Placement: here, after ledger and conditionality framing and before the eleven
% mechanism families.
% Data authority: Phase F s4_tradeoff_families.csv plus the locked sentinel
% correction. Lineage: 404 governed rows/169 studies/373 conditional rows;
% excluding two absence-status audit sentinels gives 402 substantive rows,
% 168 studies, and 371 conditional substantive rows.
% Layout: three aligned panels sharing the same eleven-family y order.
% Panel A uses stacked horizontal bars for 218 quantitative and 184 qualitative
% rows. Panel B uses lollipops for the overlapping number of studies in each
% family; the overall union is 168 and family counts must not be summed. Panel C
% plots conditional share and labels each dot conditional/rows, with 92.3%
% overall as a reference. The corrected bandwidth/resource row is 50+44=94,
% 70 studies, 89 conditional; the qualitative/partial row is 2+8=10, 8 studies,
% 7 conditional. Small n=1 families use open markers.
% Future lead-in: ``Figure~\ref{fig:tradeoff_profile} separates row type,
% overlapping study coverage, and conditionality across the eleven recurring
% coupling families.''
% Exact future caption: ``Conditionality profile for 402 substantive
% communication and sensing tradeoff rows from 168 studies. Panel (a) separates
% 218 quantitative and 184 qualitative rows. Panel (b) reports overlapping
% study counts, and panel (c) reports conditional shares within each family.
% Overall, 371 of 402 rows are conditional. The governed ledger also contains
% two absence status audit rows, bringing its total to 404. Row frequency
% describes evidence coverage, while each transfer claim remains attached to
% its source conditions.''

\subsubsection{Bandwidth, Spectrum, and Resource Allocation}

Bandwidth links the two functions when both draw on the same signal space. That
space may be divided across subcarriers, wavelengths, guard intervals,
observation time, or the usable response of a device. Embedded pilots show the
basic pattern. Increasing pilot support improves the reported pilot deviation
proxy while changing BER and the data carrying fraction
\cite{OISAC_SCR00007}. Protection intervals and chirp sweeps create related
exchanges by improving sensing isolation or pulse compression within device
bandwidth and subcarrier limits
\cite{OISAC_SCR00344,OISAC_SCR00366}. In each case, the effect follows the
resource being shared and the limit that becomes active.

Waveform structure determines how the shared signal space is used, with format
and operating mode setting time and spectral sharing
\cite{OISAC_SCR00011,OISAC_SCR00012,OISAC_SCR00083}. Within that structure,
digital designs assign subbands, chirp parameters, power, and weighted resources
\cite{OISAC_SCR00367,OISAC_SCR00626,OISAC_SCR00955,OISAC_SCR00956,OISAC_SCR01010}.
Pilot strength, placement, and synchronization fields define reference support
and communication margin, while receiver partitioning and ranging sequence
redundancy set the usable signal region
\cite{OISAC_SCR00086,OISAC_SCR00151,OISAC_SCR00369,OISAC_SCR00990,OISAC_SCR01054,OISAC_SCR01088}.
A tunable laser study supplies qualitative generation context for the same
design space \cite{OISAC_SCR00553}. Together, the allocation rule and receiver
condition define the operating point reported in each design.

Resource sharing also appears in time. Observation duration, camera exposure,
and task scheduling connect sensing quality and cadence with communication
goodput or execution time
\cite{OISAC_SCR00189,OISAC_SCR00273,OISAC_SCR00833}. Time division, chirp
repetition, coding traffic, security control, and embedded pulse count place
both functions within a shared frame
\cite{OISAC_SCR00836,OISAC_SCR01015,OISAC_SCR01034,OISAC_SCR01045,OISAC_SCR01019}.
Other designs reserve spectral, temporal, spatial, polarization, or phase
structure to manage interference, coverage, and receiver complexity
\cite{OISAC_SCR00177,OISAC_SCR00396,OISAC_SCR00916,OISAC_SCR00920,OISAC_SCR00953,OISAC_SCR00983}.
The observed outcome therefore depends on the frame period, exposure response,
isolation rule, and remaining occupancy.

When both functions occupy the same waveform, bandwidth affects pulse
compression, estimation precision, sidelobes, and radar image quality together
with communication support
\cite{OISAC_SCR00362,OISAC_SCR00884,OISAC_SCR00889,OISAC_SCR00891,OISAC_SCR00909}.
Stimulated Brillouin scattering, frequency partitioning, and constellation
choice reveal device and spectral limits
\cite{OISAC_SCR00965,OISAC_SCR00970,OISAC_SCR00988}. Doppler refinement depends
on phase noise and computation \cite{OISAC_SCR00900}. A communication and
sensing bandwidth ratio gains meaning from total bandwidth, target range, and
receiver architecture \cite{OISAC_SCR00784}. These results show that signal
bandwidth and receiver processing jointly define the operating point.

Reuse extends this mechanism by carrying sensing content in a shared waveform
or an existing transport path
\cite{OISAC_SCR00996,OISAC_SCR01004,OISAC_SCR01006}. Phase coding, carrier
grouping, radio frequency bandwidth, and space and wavelength allocation then
set the occupancy and isolation available to each function
\cite{OISAC_SCR01033,OISAC_SCR01036,OISAC_SCR01040,OISAC_SCR01086}. The result
depends on how the receiver separates the shared content and recovers the
underlying sequence.

Hardware and processing resources complete the picture. Reference recovery and
coding consume timing or redundancy
\cite{OISAC_SCR00278,OISAC_SCR00607,OISAC_SCR00730}, while emitter bandwidth,
source coherence, optical budget, and distributed learning traffic set hardware
and network capacity
\cite{OISAC_SCR00763,OISAC_SCR01071,OISAC_SCR01073,OISAC_SCR01081}. Receiver
count introduces alignment and crosstalk constraints. In the reported
prototype, four surrounding receivers were designed and one receiver was
operated \cite{OISAC_SCR00619}. Phase control and observation settings affect the
recoverable path \cite{OISAC_SCR01085,OISAC_SCR01103}, while modulation,
sideband, and time bandwidth controls affect signal recovery
\cite{OISAC_SCR01152,OISAC_SCR01153,OISAC_SCR01188}. The same nominal allocation
can therefore produce a different result when the platform or receiver changes.

\subsubsection{Power, Energy, and Dynamic Range}

Power affects both functions through the finite operating range of the
transmitter, channel, and receiver. Launch and received power, carrier and
probe ratios, device bias, modulation swing, and converter headroom refer to
different points in that chain. As power changes, performance follows the
component that first approaches a noise, nonlinear, or saturation limit. The
reference plane and limiting component are therefore part of the result.

A fiber probe sweep shows this response directly. Raising the sensing to
telecom power ratio improves sensing SNR and raises BER, with the strongest link
penalty near the probe \cite{OISAC_SCR00057}. In a visible light prototype, one
optical branch supports communication and ranging while another supports
energy harvesting. Their allocation governs rate, ranging quality, and
harvested power under separate receiver and endpoint bias conditions
\cite{OISAC_SCR01025}. Similar ratios connect BER or OSNR margin with phase,
polarization, and vibration observability in fiber
\cite{OISAC_SCR00013,OISAC_SCR00072,OISAC_SCR00220,OISAC_SCR00099,OISAC_SCR00592}.
In photonic wireless links, carrier and waveform ratios connect phase recovery,
BER or EVM, and radar or ranging quality within the available receiver headroom
\cite{OISAC_SCR00492,OISAC_SCR00954,OISAC_SCR00966,OISAC_SCR00985}.
Comparisons are most informative when they use the same power reference and
receiver margin.

Absolute power sweeps reveal where the operating range ends. In fiber, launch
and received power expose nonlinear interaction and sensing reach
\cite{OISAC_SCR00473,OISAC_SCR01079,OISAC_SCR01180}. Photonic front ends examine
operating limits associated with input power, modulator drive, converter
linearity, and amplifier gain
\cite{OISAC_SCR00631,OISAC_SCR00986,OISAC_SCR01012,OISAC_SCR01151}.
Architecture then decides where the available power is used. Polarization
paths, receiver taps, and branches assigned to individual functions distribute
power before estimation
\cite{OISAC_SCR00346,OISAC_SCR01072,OISAC_SCR01075}. Propagation, insertion,
and alignment losses determine how much reaches those branches
\cite{OISAC_SCR00002,OISAC_SCR00218,OISAC_SCR00506}. The remaining margin sets
the feasible operation of optical phased arrays, communication users, sensing
estimators, and illumination services
\cite{OISAC_SCR00940,OISAC_SCR00941,OISAC_SCR00980,OISAC_SCR01028}.

Power allocation can also be written into the waveform. Service branches,
reference copies, and relative waveform amplitudes connect radar masking,
ranging, data delivery, and energy harvesting through a shared allocation
\cite{OISAC_SCR00893,OISAC_SCR00919,OISAC_SCR01032,OISAC_SCR01165}. Sequence
transitions, coding choices, and chaotic source parameters change how
communication and sensing information is recovered
\cite{OISAC_SCR00052,OISAC_SCR00460,OISAC_SCR01105}. Frequency separation,
probabilistic shaping, sensing slots, and carrier reuse move energy and
occupancy across signal dimensions
\cite{OISAC_SCR01009,OISAC_SCR01020,OISAC_SCR01030,OISAC_SCR01093}. The
waveform definition, hardware margin, and performance threshold determine the
resulting operating point. These same factors also govern reliability and
sensing quality.

\subsubsection{Communication Reliability and Sensing Quality}

Communication reliability and sensing quality become linked after the waveform
and receiver are selected. Communication evidence uses BER, EVM, Q factor, and
required OSNR, while sensing evidence uses SNR, estimation error, and detection.
These measures can move together as leakage, nonlinear mixing, phase noise,
cancellation quality, or receiver competition changes. Some controls reduce an
interference path, while others redistribute energy within the shared signal.
This distinction explains why similar metric pairs can follow different trends.

Two studies show why interference control and power redistribution require
separate interpretation. In a fiber link combining communication with coherent
reflectometry, sequence shaping reduces nonlinear interference induced by
signal transitions. This lowers the communication penalty while retaining
phase sensing at the reported condition \cite{OISAC_SCR00052}. In a millimeter
wave waveform, DC offset strengthens the sensing component and raises radar SIR together with
communication EVM \cite{OISAC_SCR00808}. The first mechanism changes how
interference is produced, while the second redistributes energy within the
signal. Their trends therefore arise from different controls.

Architecture then determines where these effects enter the chain. Shared path
tests evaluate communication reliability beside sensing outcomes at stated
operating points and validation settings
\cite{OISAC_SCR00008,OISAC_SCR00056,OISAC_SCR00589,OISAC_SCR00680}. Dedicated
monitoring wavelengths and reused channel estimates locate sensing support in
optical filtering or receiver processing, with communication performance
described qualitatively \cite{OISAC_SCR00957,OISAC_SCR00959}. The path and
receiver therefore define the scope of each reliability result.

Once the architecture is fixed, waveform and inference choices determine the
reported communication and sensing metrics. Carrier choice, DC offset, and phase
modulation affect recovery at stated settings
\cite{OISAC_SCR00160,OISAC_SCR00312,OISAC_SCR00602}. Coding, guard structure,
and waveform composition affect BER, correlation quality, and ranging error
\cite{OISAC_SCR00885,OISAC_SCR00902,OISAC_SCR00918}. At the
inference stage, loss weighting connects communication reliability with
distance and velocity errors \cite{OISAC_SCR00233}, while Kalman tuning acts
within the sensing branch \cite{OISAC_SCR00665}. Shared sequence structure and
finite block phase processing place radar detection and phase estimation under
their sample and noise models \cite{OISAC_SCR00719,OISAC_SCR00906}. Phase index
and pulse position timing likewise connect communication error with radar SNR
or time of arrival precision \cite{OISAC_SCR00937,OISAC_SCR00979}. The sample,
noise, timing, and calibration definitions give each result its scope.

At the system level, access placement guided by channel sensing connects link
capacity with weather and demand \cite{OISAC_SCR00922}. The aggregation of
sensing traffic connects capacity with target count and resolution requirements
\cite{OISAC_SCR01049}. Simultaneous use of time and frequency
creates direct waveform coupling \cite{OISAC_SCR00978}, while shared hardware
with waveform switching creates a different path
\cite{OISAC_SCR01041}. Waveform robustness under fading and learned phase
sensing place the limiting response in the channel or processing stage
\cite{OISAC_SCR00993,OISAC_SCR01185}. These studies connect reliability and
sensing quality to the complete waveform, receiver, channel, and processing
chain. The same chain also governs rate, resolution, accuracy, and range.

\subsubsection{Rate, Resolution, Accuracy, and Range}

The link between communication rate and sensing begins with the design variable
that changes the rate. Bandwidth and symbol timing can alter resolution
directly, while frame structure, geometry, and processing change the resources
available to the estimator. The controlling variable and its operating
conditions explain why a similar rate change can produce different sensing
outcomes.

Waveform parameters provide the clearest starting point. In optoelectronic
oscillator and remote ranging designs, code length and period connect bit rate
or capacity with distance resolution and unambiguous range
\cite{OISAC_SCR00896,OISAC_SCR00576}. Power division changes OFDM reliability,
reach, and ranging error in a retroreflective prototype, while image
decomposition links stripe recovery with scene reconstruction and recognition
\cite{OISAC_SCR00912,OISAC_SCR00110}. LFM slope indexing and time division
designs in the W band place their controlling choices in slope selection and
frame duration
\cite{OISAC_SCR00361,OISAC_SCR01101}. The shared parameter sets the direction
of the relationship in each case.

Frame structure creates a second path by deciding when and where each function
receives signal support. Dynamic allocation and intensity scaling distribute
support across payload, sensing patterns, and coverage
\cite{OISAC_SCR00009,OISAC_SCR00196}. OTFS parameterization and complete grid
OFDM reuse place both functions in a common time frequency structure
\cite{OISAC_SCR00210,OISAC_SCR00274}. Separate measurement paths and mode
switching show how a similar waveform can support different forms of
integration \cite{OISAC_SCR00418,OISAC_SCR01060}. The frame and measurement
schedule are therefore part of the result.

Geometry sets the region in which those resources remain useful. Steering and
access point placement define communication and localization coverage, while
surface allocation and blockage aware access select the paths available to
each function
\cite{OISAC_SCR00036,OISAC_SCR00527,OISAC_SCR00914,OISAC_SCR01011}. Source
directivity and tag delay spacing further affect coverage, range separation,
and concurrent return recovery \cite{OISAC_SCR01024,OISAC_SCR01050}. The
reported layout is therefore part of the observed trend.

Signal construction can tie payload more closely to sensing observability.
Chip length changes occupied symbol rate together with ranging precision and
security tolerance, while embedded communication data affects distributed
sensing fidelity \cite{OISAC_SCR00387,OISAC_SCR00647}. ASK rate and guard
duration connect communication timing with range resolution and ambiguity
\cite{OISAC_SCR00887,OISAC_SCR00913}. Waveform switching provides a related
parameter coupling, with rate and resolution evaluated through separate
parameter changes \cite{OISAC_SCR00899}. In the terahertz system, the available
device band supports QPSK communication and the LFM bandwidth used to achieve
range resolution
\cite{OISAC_SCR01000}. The waveform and hardware band together define the
feasible region.

Optimization studies express these mechanisms as operating regions. Energy and
asymptotic block length define communication and detection rate pairs, while
NOMA power allocation connects aggregate rate with ranging error
\cite{OISAC_SCR00761,OISAC_SCR00925}. Other studies combine capacity with
ranging information or compare separate and joint operating configurations
\cite{OISAC_SCR00926,OISAC_SCR00944}. Probabilistic shaping and embedded data
create further regions through information rate, sidelobe level, and
estimation error \cite{OISAC_SCR00974,OISAC_SCR00975}. The selected operating
point follows the stated objective, constraints, and normalization.

Processing adds another layer because the estimator decides how much of the
available signal becomes useful sensing information. Sampling period, spectral
weighting, zero padding, and transform length exchange response time and
computation for resolution under stated acquisition settings
\cite{OISAC_SCR00238,OISAC_SCR00998,OISAC_SCR01018,OISAC_SCR01095}. Rate
selection can tune the sensitivity of a density estimate derived from BER
\cite{OISAC_SCR01021}. In a separate application, classifier size changes
recognition accuracy and edge processing burden \cite{OISAC_SCR01047}. These
studies distinguish waveform rate from the update and computation needed to
interpret it.

At network scale, the relationship moves from a single waveform to the traffic
and spectrum available across the system. Transport capacity controls how many
sensing streams meet an echo quality threshold \cite{OISAC_SCR00436}.
Separated bands, idle time frequency cells, cyclic prefix correlation, and
shared optical waveforms connect occupied spectrum with resolution or
estimation support
\cite{OISAC_SCR00977,OISAC_SCR00981,OISAC_SCR01023,OISAC_SCR01115}. A visible
light study carries this logic into a proposed capacity and localization
problem \cite{OISAC_SCR00488}. The evidence thus distinguishes observed
operation from a proposed design space.

\subsubsection{Waveform, Hardware, and Complexity}

Implementation complexity depends on where a design places its burden. Source
coherence, receiver structure, calibration, synchronization, computation,
memory, and latency represent different costs. Reusing a component may move
work into processing or control, so implementation complexity must be assessed
across the complete chain.

Two photonic systems show how this transfer occurs. A terahertz design uses
coherent comb generation to stabilize the QPSK and LFM paths and concentrates
complexity in the optical source chain \cite{OISAC_SCR01000}. A long span fiber
experiment uses lower cost lasers with broad linewidths and recovers coherent
data and vibration through residual carrier processing
\cite{OISAC_SCR01070}. Each architecture preserves performance through a
different balance of source hardware and digital processing.

\subsubsection{Other Implementation Relationships}

Other implementation relationships show how a shared choice appears at
different layers of the system. Its effect may emerge as protection, joint
gain, device sensitivity, estimator burden, control freshness, or allocation
latency. Each case carries its own conditions and prepares the comparison
limits that follow.

Security and synergy illustrate two possible directions. Phase keyed recovery
uses ranging tolerance to create a physically bounded decryption region
\cite{OISAC_SCR00387}. With spatial alignment, communication gain and target
profile estimation can benefit from the same optical design
\cite{OISAC_SCR00941}. Their geometry, objective, and validation setting explain
why one relationship expresses protection and the other expresses joint gain.

Other cases locate the relationship within the implementation chain. Electrode
transparency connects device architecture with photodetector sensitivity, while
passive return loss and repeater structure affect access for fault localization
\cite{OISAC_SCR00001,OISAC_SCR01100}. Capacity and distance estimation define a
model operating region \cite{OISAC_SCR00908}. Queue freshness enters the robotic
control loop, and local spectral scanning connects computation with range
resolution \cite{OISAC_SCR00910,OISAC_SCR00994}. Weighted priorities guide
resource allocation, while orthogonal task channels shape application latency
\cite{OISAC_SCR01001,OISAC_SCR01059}. Read together, these
cases locate the relationship at the device, estimator, control, or allocation
layer and provide the conditions needed for comparison across studies.

\subsection{Conditions for Comparison Across Studies}

Comparison across studies begins with alignment. The bounded relations summarized
in Table~\ref{tab:metric_admissibility} connect studies when their task, metric
definition, measurement plane, operating point, and validation setting can be
matched through the framework in Section~\ref{sec:foundations_comparison}.
Other metric evidence explains mechanisms, feasibility, and behavior within
individual systems. This separation keeps each reported value attached to its
original engineering meaning.

Tradeoff evidence follows the same principle. Because objectives and
constraints vary across systems, the synthesis compares the shared mechanism
and the conditions that produce the observed response. This approach shows why
a relationship recurs and when it can guide another design. Its engineering
value also depends on the validation setting and on how completely the
configuration can be reconstructed. Section~\ref{sec:validation_reproducibility}
examines these requirements through validation settings, field evidence across
both functional domains, and reconstructability.
